\chapter{Fundamentos de la supervisión y protección infantil digital}
\label{chap:fundamentos-supervision}

\insertminitoc
\parindent1.5em

\section{Supervisión parental y dinámicas de mediación en el hogar moderno}

La supervisión parental en la era de la movilidad inteligente ha evolucionado desde una vigilancia física directa hacia una gestión abstracta de riesgos digitales invisibles. Según las investigaciones de \cite{ayyash_cybersecurity_2024}, aunque el 87\% de los padres considera la ciberseguridad una prioridad máxima para su hogar, existe una parálisis operativa debido a que las herramientas actuales suelen ser demasiado complejas o intrusivas. Esta sección analiza cómo la mediación parental debe equilibrar la necesidad de protección con el derecho fundamental del menor a la privacidad y autonomía. La literatura contemporánea distingue claramente entre la mediación restrictiva, que se basa en la prohibición y el filtrado, y la mediación habilitadora, que utiliza la tecnología como un puente para fomentar el uso seguro a través del diálogo y la capacitación del menor \cite{livingstone_maximizing_2017}. El desafío reside en que la supervisión técnica a menudo se percibe como una ruptura de la confianza, lo que puede llevar a los menores a ocultar sus actividades en lugar de reportar riesgos reales.

Investigaciones recientes sugieren que las herramientas de supervisión son más efectivas cuando actúan como un sistema de "visibilidad contextualizada", permitiendo que el tutor intervenga solo ante incidentes críticos detectados mediante análisis inteligente \cite{stoilova_parental_2024}. Un factor distorsionador en esta dinámica es la "tecnoferencia", un fenómeno donde la propia dependencia de los padres hacia sus dispositivos interrumpe la supervisión activa y reduce la calidad de la interacción familiar \cite{selak_problematic_2025}. Esta falta de atención parental empuja a menudo a los menores a buscar validación en redes sociales, aumentando su exposición a riesgos externos. Por ello, la mediación debe ser un proceso bidireccional; estudios de co-diseño con menores revelan que estos aceptan la presencia de software de seguridad si comprenden que su fin es pedagógico y no simplemente punitivo \cite{mcnally_co-designing_2018}.

Además, la efectividad de la mediación está intrínsecamente ligada al desarrollo de habilidades digitales tanto en padres como en hijos. No basta con instalar un filtro; es necesario que el sistema proporcione "andamiaje" educativo que desaparezca gradualmente a medida que el niño madura y demuestra capacidad de autorregulación. Los modelos de supervisión modernos deben evitar el "monitoreo ciego" y favorecer una transparencia algorítmica donde el menor sepa qué se está supervisando y por qué. Este enfoque reduce la reactancia psicológica, un estado donde el niño intenta eludir activamente las restricciones solo por el hecho de sentirlas arbitrarias \cite{ayyash_cybersecurity_2024}. En última instancia, la tecnología debe servir para fortalecer el vínculo familiar, convirtiendo cada incidente digital en una oportunidad de aprendizaje compartido que prepare al menor para una ciudadanía digital autónoma.


\section{Prevención de riesgos, victimización y salud pública digital}

La protección de menores en el entorno digital ha trascendido el ámbito de la seguridad informática para convertirse en una prioridad de salud pública global. La exposición a contenidos inapropiados, el ciberacoso y la victimización sexual en línea (OCSEA) generan secuelas psicológicas que pueden ser tan devastadoras como el abuso físico, afectando la resiliencia y el desarrollo emocional a largo plazo \cite{page_psychological_2025}. Esta sección explora la epidemiología de estos riesgos, destacando que la detección temprana es la única vía para mitigar daños acumulativos. Datos recientes indican que casi el 17.4\% de los menores que navegan de forma autónoma sufren una exposición involuntaria a la pornografía, a menudo impulsada por fallos en los algoritmos de recomendación de las plataformas comerciales \cite{yusuf_lets_2023}. Esta realidad exige sistemas inteligentes capaces de identificar no solo el contenido explícito, sino también las semánticas de riesgo antes de que se produzca el contacto perjudicial.

La vulnerabilidad psicológica actúa como un catalizador en los procesos de victimización. Los menores que ya presentan niveles elevados de distrés emocional o baja autoestima son más susceptibles a las tácticas de manipulación de los depredadores en línea, quienes explotan su necesidad de pertenencia \cite{el-asam_psychological_2022}. Esta victimización suele escalar rápidamente en entornos cifrados, donde los métodos tradicionales de monitoreo de red no tienen visibilidad. Además de los riesgos externos, el propio comportamiento sedentario y el uso excesivo de pantallas durante la noche se han correlacionado directamente con el desarrollo de síntomas depresivos y ansiedad en niños y adolescentes \cite{soriano-molina_association_2025}. Por tanto, una prevención integral debe contemplar tanto la seguridad de los contenidos como la higiene digital del menor, alertando sobre patrones de uso que comprometan su bienestar físico y mental.

Un modelo preventivo robusto debe actuar como un "escudo cognitivo", filtrando influencias dañinas mientras fomenta la capacidad crítica del usuario. La literatura subraya que el ciberacoso es uno de los riesgos más prevalentes, con tasas de victimización que alcanzan el 57.5\% en algunos contextos globales, siendo la violencia verbal el tipo más común \cite{zhu_cyberbullying_2021}. Las herramientas de supervisión deben, por tanto, ser capaces de detectar el tono y el sentimiento de las conversaciones en tiempo real, proporcionando al padre una bitácora de evidencias que sirva para intervenir antes de que el acoso derive en crisis graves de salud mental. Al integrar la detección visual y lingüística, es posible crear un entorno controlado donde el menor pueda explorar las ventajas educativas de la red sin quedar desprotegido ante las dinámicas hostiles que caracterizan gran parte de la web moderna \cite{youssef_online_2025}.

\section{Arquitecturas de monitoreo y limitaciones de los sistemas comerciales}

A pesar de la proliferación de aplicaciones de control parental en el mercado, la mayoría presenta vulnerabilidades técnicas críticas que comprometen su efectividad real. Investigaciones profundas de \cite{rojas_technical_2025} han demostrado que los métodos tradicionales basados en filtrado de DNS y listas negras de URLs son fácilmente eludidos por menores con conocimientos básicos, mediante el uso de VPNs o navegadores que no respetan las políticas del sistema operativo. Esta sección analiza por qué la supervisión debe desplazarse desde la capa de red hacia la capa de interfaz de usuario (HCI). Al realizar capturas de pantalla y procesarlas localmente, se obtiene una visión real de lo que el niño consume, superando las limitaciones de los logs de navegación que no capturan la actividad dentro de aplicaciones de mensajería cerrada o redes sociales de video efímero \cite{koulaxidis_improving_2022}.

Sin embargo, el monitoreo técnico enfrenta un dilema ético y de seguridad masivo relacionado con la privacidad de los datos. Un análisis exhaustivo de \cite{feal_angel_2020} reveló que muchas de las aplicaciones comerciales más populares abusan de los servicios de accesibilidad de Android para recolectar datos sensibles, que luego son vendidos a terceros para fines publicitarios. Este hallazgo convierte a la propia herramienta de protección en un riesgo de privacidad para la familia. El modelo propuesto en esta investigación rompe con esa tendencia al priorizar el procesamiento local (*on-device*), garantizando que las evidencias capturadas permanezcan cifradas en el dispositivo del tutor y nunca toquen servidores externos. Este enfoque no solo blinda la privacidad, sino que también evita que la aplicación se convierta en un vector de ataque para ciberdelincuentes interesados en los perfiles de comportamiento infantil.

Otro desafío técnico crítico es la gestión de la "fatiga de alerta". Los sistemas basados en reglas rígidas suelen generar una tasa de falsos positivos inasumible para los padres, quienes terminan desactivando las notificaciones y dejando al menor desprotegido \cite{rojas_technical_2025}. La solución reside en la implementación de una Inteligencia Artificial capaz de discernir el contexto; por ejemplo, diferenciar una imagen médica de una imagen pornográfica, o un lenguaje coloquial entre amigos de un acoso real. Al optimizar la precisión de estas alertas, se asegura que el tutor solo reciba notificaciones ante eventos de riesgo genuino, manteniendo la relevancia de la herramienta. La arquitectura debe ser ligera para no degradar el rendimiento del sistema operativo del menor, evitando cierres inesperados por parte del gestor de recursos de Android que dejarían al sistema inoperativo en momentos críticos.


\section{Inteligencia Artificial ética y transparencia en el diseño para menores}

La integración de la Inteligencia Artificial (\gls{ia}) en sistemas de protección infantil introduce dilemas éticos que deben abordarse desde el diseño inicial. El uso de algoritmos predictivos para evaluar riesgos conlleva el peligro de sesgos algorítmicos que pueden discriminar basándose en jerga cultural, dialectos regionales o el origen socioeconómico del menor \cite{agbana_ethical_2025}. Esta sección examina la necesidad de una \gls{ia} explicable (*Explainable AI*), donde el sistema no solo bloquee un contenido, sino que sea capaz de justificar ante el tutor la razón técnica y ética detrás de dicha decisión. La transparencia algorítmica es vital para que el padre pueda validar la precisión del sistema y para que el menor aprenda a identificar por sí mismo los criterios de peligrosidad en la red \cite{balapour_mobile_2020}.\\

La ética del monitoreo también exige considerar la huella digital que el sistema genera sobre el menor. La recolección masiva de datos desde la primera infancia crea un perfil digital que puede condicionar el futuro académico o profesional del individuo si no se gestiona con principios de minimización de datos \cite{swider-cios_young_2023}. Por esta razón, los estándares internacionales propuestos por la OECD \cite{noauthor_full_2025} subrayan que las herramientas de bienestar digital deben ser "seguras por defecto" y diseñadas para proteger la identidad del niño frente a la comercialización. El modelo de bitácora propuesto en esta tesis se alinea con estos principios al permitir que el tutor sea el único poseedor de la clave de descifrado de las evidencias, asegurando que la vigilancia sea proporcional y necesaria según el marco legal vigente (\gls{gdpr} y COPPA).\\

Además, es imperativo mantener al "humano en el ciclo" (*human-in-the-loop*). La \gls{ia} no debe tomar decisiones punitivas autónomas, sino actuar como un sensor avanzado que empodera al tutor para tomar la decisión pedagógica final. Según \cite{agbana_ethical_2025}, la dependencia excesiva de la automatización puede erosionar la intuición parental y la comunicación intrafamiliar. La herramienta inteligente debe ser vista como un copiloto que reduce el ruido informativo, pero que nunca sustituye el juicio ético del padre. Al involucrar a los menores en el entendimiento de cómo funcionan estos algoritmos, se fomenta un pensamiento crítico esencial para navegar en un mundo donde la \gls{ia} mediará casi todas sus interacciones futuras, promoviendo una resiliencia digital que trasciende la simple restricción técnica.\\

\section{Optimización para el procesamiento en el borde (Edge AI) y rendimiento}

Para que una solución de supervisión inteligente sea viable en dispositivos móviles con recursos limitados, es fundamental optimizar el procesamiento en el borde, conocido como *Edge AI*. Esta sección detalla las ventajas técnicas de ejecutar modelos de aprendizaje profundo directamente en el hardware del teléfono, evitando la latencia y los riesgos de privacidad asociados con la nube \cite{sahu_edge_2025}. La implementación de arquitecturas ligeras permite realizar análisis de imágenes y texto en milisegundos sin comprometer la fluidez de las aplicaciones que el menor está utilizando, un factor crítico para evitar que el niño desinstale la herramienta por razones de rendimiento \cite{ali_efficient_2021}. El uso de técnicas como la cuantización de modelos y la destilación de conocimiento permite que redes neuronales complejas se ejecuten en procesadores móviles de gama media sin agotar excesivamente la batería.\\

Un aspecto técnico a menudo ignorado es la interacción entre el software de monitoreo y el kernel del sistema operativo. Investigaciones de \cite{cao_mobile_2023} sugieren que la gestión eficiente de la memoria y la rectificación del reclamo de páginas son vitales para aplicaciones que deben permanecer activas en segundo plano. El uso inteligente de *zram* (memoria comprimida en \gls{ram}) en dispositivos Android permite que la aplicación de supervisión mantenga su estado sin causar "thrashing" o ralentizaciones en el cambio de aplicaciones del menor. Esta optimización de bajo nivel es la que garantiza que la protección sea continua y no se vea interrumpida por el propio sistema operativo en momentos de alta carga de trabajo. Al procesar las evidencias localmente, se elimina también el costo de datos móviles que supondría subir capturas de pantalla de alta resolución a un servidor, haciendo la solución más accesible para familias con planes de datos limitados.\\

Finalmente, la robustez técnica del sistema debe estar orientada a fomentar la resiliencia digital a largo plazo. El objetivo último de la tecnología propuesta es su propia obsolescencia; es decir, que el menor desarrolle las funciones ejecutivas necesarias para navegar de forma segura sin asistencia técnica \cite{clemente-suarez_digital_2024}. La bitácora de evidencias sirve aquí como una herramienta de retroalimentación en tiempo real, donde el procesamiento en el borde permite dar consejos instantáneos al menor en el momento exacto en que ocurre un riesgo. Esta inmediatez pedagógica, apoyada por una infraestructura técnica eficiente y privada, sienta las bases para un ecosistema digital donde la tecnología actúa como un andamiaje protector que empodera el desarrollo cognitivo y emocional del niño en un mundo irreversiblemente conectado \cite{holmarsdottir_integrative_2025, sahu_edge_2025}.\\